\documentclass[11pt]{article}
\usepackage[a4paper,margin=1in]{geometry}
\usepackage{amsmath,amssymb,mathtools,bm}
\usepackage{enumitem,booktabs,array}
\usepackage{tcolorbox}
\usepackage{hyperref}
\usepackage[protrusion=true,expansion=false]{microtype}
\hypersetup{colorlinks=true,linkcolor=blue,urlcolor=blue,citecolor=blue}
\setlist[itemize]{topsep=2pt,itemsep=2pt}
\setlist[enumerate]{topsep=2pt,itemsep=2pt}
\newtcolorbox{keybox}{colback=blue!3!white,colframe=blue!40!black,boxrule=0.4pt,arc=1mm}
\newtcolorbox{alertbox}{colback=red!3!white,colframe=red!40!black,boxrule=0.4pt,arc=1mm}
\newtcolorbox{answerbox}{colback=green!3!white,colframe=green!40!black,boxrule=0.4pt,arc=1mm}
\newtcolorbox{drillbox}{colback=yellow!5!white,colframe=orange!60!black,boxrule=0.4pt,arc=1mm}
\newcommand{\EV}{\mathbb{E}}
\newcommand{\Var}{\mathrm{Var}}
\newcommand{\Cov}{\mathrm{Cov}}
\newcommand{\blank}[1]{\underline{\hspace{#1}}}

\title{W05-04: Jiang, Matvos, Piskorski \& Seru (2024) \\
\large ``Monetary Tightening and U.S. Bank Fragility in 2023: Mark-to-Market Losses and Uninsured Depositor Runs?'' --- Active Recall Workbook}
\author{Banking \& Risk Management --- Exam Prep}
\date{\today}
\begin{document}\maketitle

\begin{keybox}
\textbf{Paper in one sentence.} JMPS show that Fed rate hikes 2022--23 generated $\$2.2$ trillion of mark-to-market losses on US bank asset portfolios (about $10\%$ of assets); if half of uninsured depositors run at banks whose MTM-adjusted equity is most impaired, $\sim 186$ banks with $\$300\text{B}+$ of uninsured deposits are at risk of run-induced insolvency --- the quantitative frame behind the Silicon Valley Bank, Signature and First Republic failures.

\textbf{Placement.} The Silicon-Valley-Bank post-mortem paper; modern counterpart to Gorton--Metrick and Egan--Hortaçsu--Matvos on run dynamics, extended to rising-rate-driven fragility.
\end{keybox}

\section*{Round 1: Complete Walk-Through}

\subsection*{1.1 The macro setup}
March 2022--early 2023: Federal funds rate rises from $\sim 0\%$ to $\sim 5\%$. Long-duration bank assets (treasuries, MBS) lose value. Under GAAP, held-to-maturity (HTM) securities can be carried at amortised cost, \emph{hiding} unrealised losses from book equity. JMPS reprice bank assets to market.

\subsection*{1.2 Central calculation}
For each US bank (Call Reports), JMPS:
\begin{enumerate}
\item Estimate asset-duration using reported maturity buckets.
\item Apply the rate-hike-induced price change to each bucket.
\item Sum unrealised losses to get aggregate MTM losses.
\end{enumerate}
Aggregate MTM loss on US bank assets: $\sim \$2.2$T, or $10\%$ of total assets.

\subsection*{1.3 Uninsured-deposit run vulnerability}
Define a bank as \emph{run-at-risk} if MTM-adjusted equity $<$ losses that would occur if a fraction $\theta$ of uninsured depositors run and the bank must fire-sell. For $\theta=50\%$:
\begin{itemize}
\item $\sim 186$ US banks have insufficient MTM equity to withstand the run.
\item They hold roughly $\$300$B of uninsured deposits.
\item Silicon Valley Bank was a textbook case: long-duration HTM, concentrated uninsured depositor base (VC-backed tech startups).
\end{itemize}

\subsection*{1.4 Identification of run risk}
\begin{itemize}
\item Cross-section of banks with varying asset-duration and uninsured-deposit share.
\item Banks with large MTM losses and large uninsured-deposit shares are most fragile.
\item Simple composite index: $\text{Fragility} = \text{MTM loss} / \text{Equity} + \text{UninsDep} / \text{Liab}$.
\end{itemize}

\subsection*{1.5 Policy implications}
\begin{itemize}
\item HTM accounting hides duration risk from supervisors.
\item Expansions of deposit insurance (FDIC's 2023 systemic-risk exception for SVB) internalise run externalities.
\item Potential liquidity rules: LCR / HQLA requirements do not fully capture duration risk at the asset side.
\end{itemize}

\subsection*{1.6 Caveats}
\begin{alertbox}
\begin{itemize}
\item Static mark-to-market: does not model deposit-side repricing (which moderates equity-value loss) --- DSS deposits-channel correction.
\item Run parameter $\theta$ is a choice; headline result is robust across 25--75\%.
\item Cross-section rather than time series; counterfactual ``what if rates hadn't risen'' is implicit.
\end{itemize}
\end{alertbox}

\section*{Round 2: Level 1 Fill-in}
\begin{enumerate}
\item Macro shock: Federal funds rate rises $\sim$\blank{1cm}\% $\to$ $\sim$\blank{1cm}\%.
\item Aggregate MTM loss: $\sim$\$\blank{1cm} trillion, $\sim$\blank{1cm}\% of assets.
\item Number of banks at run-risk (50\% run): $\sim$\blank{1cm}.
\item Uninsured deposits at those banks: $\sim$\$\blank{1cm} billion.
\item Headline case study: \blank{2cm} Valley Bank.
\end{enumerate}

\section*{Round 3: Level 2 Prompts}
\begin{enumerate}
\item Derive the duration-loss formula: $\Delta P/P \approx -D\cdot\Delta y$. How does this enter JMPS?
\item Why does HTM accounting hide fragility?
\item Write the run-threshold condition: MTM equity vs.\ fire-sale loss.
\item Compare JMPS's mechanism to Gorton--Metrick on repo runs.
\end{enumerate}

\section*{Round 4: Minimal Scaffolding}
From a blank page: (i) rate shock, (ii) MTM revaluation, (iii) run-threshold calculation, (iv) number of at-risk banks, (v) two caveats.

\section*{Round 5: Blank Page}
Essay: the 2023 regional-bank crisis as a duration-driven run --- what JMPS teach us.

\vspace{6pt}
\begin{drillbox}
\textbf{Speed Drill.} (1) Rate shock range. (2) Aggregate MTM loss. (3) Fragility count. (4) Canonical failure case. (5) Role of HTM accounting.
\end{drillbox}

\section*{Quick Reference \& Answer Key}
\begin{answerbox}
\begin{itemize}
\item \textbf{Shock.} Fed hikes $\sim 0\%\to\sim 5\%$, 2022--23.
\item \textbf{MTM.} $\sim\$2.2$T loss, $\sim 10\%$ of bank assets.
\item \textbf{Fragility.} $\sim 186$ banks, $\sim\$300$B uninsured deposits at risk.
\item \textbf{Case.} Silicon Valley Bank --- long duration HTM $+$ concentrated uninsured base.
\end{itemize}
\end{answerbox}

\begin{keybox}
\textbf{30-second exam pitch.} Jiang, Matvos, Piskorski \& Seru (2024) produce the canonical post-mortem on the 2023 US regional banking crisis. The Fed's 2022--23 rate hikes created roughly \$2.2 trillion of unrealised mark-to-market losses on US bank asset portfolios, about 10\% of total assets. Because held-to-maturity (HTM) securities are carried at amortised cost, book equity masked the solvency hit. JMPS combine MTM losses with uninsured-deposit shares to identify banks vulnerable to a 50\%-run scenario: $\sim 186$ US banks with $\sim\$300$B of uninsured deposits sit on the fragile frontier. Silicon Valley Bank was the textbook case --- long-duration Treasuries and MBS held-to-maturity financed by a concentrated uninsured deposit base. The paper quantifies why duration risk, HTM accounting and uninsured-deposit concentration jointly drove the crisis, and motivates expansions of deposit insurance and better duration-risk supervision.
\end{keybox}

\end{document}
